\subsection{开映射定理}\label{sec:open-mapping-theorem}

设 $X$ 和 $Y$ 为 Banach 空间. 以下关于从 $X$ 到 $Y$ 的算子的三个定理, 对算子的求逆具有重要意义.

\begin{Theorem}{开映射定理}
\label{thm:open-mapping}
设 $A\in\mathcal{L}(X,Y)$, 并假设 $A$ 为满射 (映满). 那么每个开集 $S\subset X$ 的像在 $Y$ 中都是开集.
\end{Theorem}

\begin{Theorem}{有界逆定理}
\label{thm:bounded-inverse}
设 $A\in\mathcal{L}(X,Y)$, 并假设 $A$ 为双射 (一一且映满). 那么逆映射 $A^{-1}$ 有界.
\end{Theorem}

\begin{Theorem}{闭图定理}
\label{thm:closed-graph}
设 $(\mathcal{D}(A),A)$ 是从 $X$ 到 $Y$ 的线性算子. 若 $(\mathcal{D}(A),A)$ 是闭算子, 且其定义域 $\mathcal{D}(A)$ 在 $X$ 中是闭的, 则该算子有界.
\end{Theorem}

通常的泛函分析课程利用 Baire 纲定理证明开映射定理, 然后将有界逆定理和闭图定理作为推论导出. 事实上, 正如下面的 Lemma~\ref{lem:three-operator-theorems-equivalence} 将表明的, 三个定理彼此等价. 此外, 在 Section~\ref{sec:adjoint-operators} 中建立伴随算子的工具后, 就能够对 Hilbert 空间证明有界逆定理. 由于大多数应用都将使用 Hilbert 空间, 这里只证明这一情形, 更一般情形的证明请读者参阅文献.

\begin{Lemma}
\label{lem:three-operator-theorems-equivalence}
开映射定理, 有界逆定理和闭图定理彼此等价.
\end{Lemma}

这里给出的证明仅在 $X$ 是 Hilbert 空间时有效. 不过, 证明的前三部分直接适用于 Banach 空间, 最后一部分可以通过使用等价类来推广. (这是绕过 Banach 空间中缺少投影定理这一困难的标准技巧, 但本书不展开讨论.)

\begingroup
\postdisplaypenalty=10000
\renewcommand{\endProof}{\par\nobreak\hfill$\blacksquare$\par\vspace{0.1cm}}
\begin{Proof}
\textbf{开映射定理 $\Rightarrow$ 有界逆定理.} 由有界逆定理的假设立即可知, 存在定义域为 $\mathcal{R}(A)$ 的线性逆算子 $A^{-1}$. 非平凡的断言在于 $A^{-1}$ 有界. 然而, 这来自开映射定理, 线性算子的有界性与连续性的等价性, 以及连续性定义的拓扑表述: 算子 $T$ 连续, 当且仅当 $\mathcal{R}(T)$ 中开集的原像在 $\mathcal{D}(T)$ 中是开的. ($\mathcal{R}(A^{-1})$ $(=X=\mathcal{D}(A))$ 中开集在 $A^{-1}$ 下的原像, 与该集合在 $A$ 下的像相同.)

\textbf{有界逆定理 $\Rightarrow$ 闭图定理.} 首先观察到, 乘积空间 $X\times Y$ 是 Banach 空间, 其范数为
\begin{equation}\label{eq:banach-product-sum-norm}
\|(x,y)\|=\|x\|+\|y\|.
\end{equation}
假设 $\Gamma(A)$ 是 $X\times Y$ 中的闭子空间, 而 $\mathcal{D}(A)$ 是 $X$ 中的闭子空间. 因而 $\Gamma(A)$ 和 $\mathcal{D}(A)$ 都是 Banach 空间. 现在定义投影映射
\begin{equation}\label{eq:graph-domain-projection-map}
P:\Gamma(A)\to\mathcal{D}(A)
\end{equation}
为
\begin{equation}\label{eq:graph-domain-projection-action}
P(x,Ax):=x.
\end{equation}
注意, $P$ 为线性双射. 事实上, 它的逆
\begin{equation}\label{eq:graph-domain-projection-inverse-map}
P^{-1}:\mathcal{D}(A)\to\Gamma(A)
\end{equation}
定义为
\begin{equation}\label{eq:graph-domain-projection-inverse-action}
P^{-1}x:=(x,Ax).
\end{equation}
映射 $P$ 也有界, 因为
\begin{equation}\label{eq:graph-domain-projection-boundedness}
\|P(x,Ax)\|=\|x\|\leq\|x\|+\|Ax\|=\|(x,Ax)\|.
\end{equation}
因而由有界逆定理 (Theorem~\ref{thm:bounded-inverse}), 存在常数 $C$, 使得
\begin{equation}\label{eq:bounded-graph-projection-inverse-estimate}
\|(x,Ax)\|=\|P^{-1}x\|\leq C\|x\|.
\end{equation}
但这说明 $A$ 有界, 因为
\begin{equation}\label{eq:closed-graph-operator-boundedness-estimate}
\|Ax\|\leq\|(x,Ax)\|\leq C\|x\|
\end{equation}
对每个 $x\in\mathcal{D}(A)$ 成立.

\textbf{闭图定理 $\Rightarrow$ 有界逆定理.} 这一部分留作习题. (Problem~\ref{ex:closed-graph-implies-bounded-inverse}.)

\textbf{有界逆定理 $\Rightarrow$ 开映射定理.} 只在 $X$ 是 Hilbert 空间的情形下证明. 由于 $A$ 有界, $\mathcal{N}(A)$ 是闭的 (参见 Problem~\ref{ex:closed-operator-null-space-closed}). 因而可以利用投影定理将 $X$ 分解为 $X=\mathcal{N}(A)\oplus\mathcal{N}(A)^\perp$. 然后令 $P:X\to\mathcal{N}(A)^\perp$ 为正交投影算子, 并将 $\widetilde A$ 定义为 $A$ 在定义域 $\mathcal{N}(A)^\perp$ 上的限制. 注意, $A$ 可写成这两个算子的复合; 即\[Ax=\widetilde A(Px)\]对每个 $x\in X$ 成立. 现在证明的关键在于以下两个事实, 请读者验证.\begin{enumerate}\item 投影映射 $P$ 将 $X$ 中的开集映为 $\mathcal{N}(A)^\perp$ 中的开集 (Problem~\ref{ex:orthogonal-projection-open-map}).\item 算子 $\widetilde A$ 是从 $\mathcal{N}(A)^\perp$ 到 $Y$ 的连续双射 (Problem~\ref{ex:surjective-hilbert-operator-null-complement-bijection}).\end{enumerate}现在, $X$ 中的开集经 $P$ 映为 $\mathcal{N}(A)^\perp$ 中的开集, 而由有界逆定理, 这个集合经 $\widetilde A$ 映为 $Y$ 中的开集. (集合在 $\widetilde A$ 下的像是集合在 $\widetilde A^{-1}$ 下的原像.) 因而作为这两个映射的复合的映射 $A$, 将开集映为开集.
\end{Proof}
\endgroup
